\documentclass[12pt, a4paper]{article}
\usepackage{amsmath, amsthm, amssymb, mathtools}
\usepackage[margin=1in]{geometry}
\usepackage{hyperref, cleveref}
\usepackage{graphicx, subcaption}
\usepackage{natbib}
\usepackage{booktabs}

\newtheorem{theorem}{Theorem}
\newtheorem{lemma}[theorem]{Lemma}
\newtheorem{proposition}[theorem]{Proposition}
\newtheorem{conjecture}[theorem]{Conjecture}
\newtheorem{definition}{Definition}
\newtheorem{remark}{Remark}

% Key notation
\newcommand{\Cstar}{\mathbb{C}^*}
\newcommand{\lam}{\lambda}
\newcommand{\rhostar}{\rho_*}
\newcommand{\psihat}{\hat{\psi}}
\newcommand{\dd}{\mathrm{d}}

\title{Why Rotating Fluids Make Polygons:\\
       Onsager Condensation on the M\"obius Parameter Space}
\author{}
\date{}

\begin{document}
\maketitle

\begin{abstract}
Saturn has a hexagon.  Jupiter has an octagon and a pentagon.  Why do
rotating planetary atmospheres produce regular polygons?

We propose an answer through a framework built on five results.
The M\"obius conjugacy theorem places circular orbits at the locus
$|\lam| = 1$ in the complex multiplicative group $\Cstar$;
$\mathbb{Z}_N$ symmetry from the azimuthal wavenumber converts these
circles into $N$-gon jet meanders.  Within the WKB approximation, the
eigenvalue branch structure of the Rayleigh--Kuo equation disconnects
the circular ($\sigma = 0$) and radial ($\alpha = 0$) states, preventing
smooth perturbations from destroying a polygon.  Among $N$ point vortices
on a ring, the regular $N$-gon maximizes the interaction energy along
the loxodromic deformation direction (verified computationally for
$N = 3, \ldots, 9$)---the signature of negative temperature in Onsager's
(1949) two-dimensional vortex statistics.  At negative temperature, the
polygon is the most probable macrostate; atmospheric turbulence provides
the thermostat.

Three planetary configurations are tested: Saturn's hexagon ($N = 6$),
Jupiter's north polar octagon ($N = 8$, $\sigma_{\mathrm{geom}} = 0.06$),
and south polar pentagon ($N = 5$, $\sigma_{\mathrm{geom}} = 0.05$).
Jupiter's cyclone rings sit close to $|\lam| = 1$; Saturn's hexagon,
with $\sigma_{\mathrm{geom}} = 0.50$ ($r = 1.65$), departs further,
reflecting its larger meander amplitude.  All three arise through
different dynamical mechanisms---Rossby stationarity for Saturn, vortex
self-organization for Jupiter---yet converge toward the same geometric
locus.  \emph{Which} polygon forms depends on planetary parameters;
\emph{that} a polygon forms is universal.
\end{abstract}

\section{Introduction}
\label{sec:introduction}

In 1988, David Godfrey examined Voyager images of Saturn's north pole and
found a hexagon \citep{Godfrey1988}.  Not an approximate hexagon, not a
roughly six-sided shape, but a persistent, regular, six-sided polygon
inscribed in the jet stream at $76^\circ$\,N---stable across decades,
seasons, and the passage of the Cassini mission through 2014.  Thirty
years later, the Juno spacecraft revealed that Jupiter's poles harbor
their own polygonal arrangements: eight cyclones in an octagonal ring at
the north pole, five in a pentagonal ring at the south
\citep{Adriani2018}.  Laboratory rotating tanks produce similar patterns
\citep{BarbosaAguiar2010}.

These observations raise a question that is simple to state and
surprisingly difficult to answer: \emph{why do rotating fluids make
polygons?}

The question has two parts.  The first---\emph{how} does a specific
polygon form?---has received considerable attention.  Numerical
simulations reproduce Saturn's hexagon from jet instabilities
\citep{Morales2011}; vortex dynamics models capture Jupiter's cyclone
rings; Rossby wave theory explains the wavenumber selection $n = 6$.
The second part---\emph{why} should different physical mechanisms on
different planets all produce the same geometric outcome?---has not
been addressed.

This paper answers both parts.  The answer has five steps:
\begin{enumerate}
  \item \textbf{Geometry.}  The M\"obius conjugacy theorem places
    \emph{circular} orbits at the locus $|\lam| = 1$ in the complex
    multiplicative group $\Cstar$, where logarithmic spirals transition
    to circles.  A separate geometric step (Proposition~\ref{prop:polygon})
    shows that imposing $\mathbb{Z}_N$ azimuthal symmetry converts these
    circles into $N$-gon jet meanders.

  \item \textbf{Branch disconnection.}  Within the WKB approximation,
    the Rayleigh--Kuo eigenvalue equation has two branches---elliptic
    ($\sigma = 0$, polygonal) and hyperbolic ($\alpha = 0$, radial)---that
    are disconnected.  Smooth perturbations of the jet speed cannot move
    the system between branches at leading order.

  \item \textbf{Energy maximum.}  Among $N$ point vortices on a ring, the
    regular $N$-gon maximizes the interaction energy: $H''(0) < 0$ for
    all $N \ge 3$ (proven via a closed-form formula that is manifestly
    negative).

  \item \textbf{Onsager condensation.}  In Onsager's (1949) statistical
    mechanics of two-dimensional vortices, systems at negative temperature
    are driven toward energy maxima.  The polygon is therefore the most
    probable macrostate in the large-scale limit of 2D turbulence.

  \item \textbf{Atmospheric thermostat.}  Small-scale turbulence in
    planetary atmospheres provides the effective mixing that enables
    Onsager relaxation.  The polygon is the endpoint.
\end{enumerate}
Steps 1--3 are proven within this paper (Step~3 via an explicit
closed-form formula for $H''(0)$ that is manifestly negative for all
$N \ge 3$).  Step 4 is established statistical mechanics.  Step 5 is a
physical observation.

The framework is universal: it does not depend on the specific dynamical
mechanism that produces the polygon.  Saturn's hexagon arises from Rossby
wave stationarity (a jet meander); Jupiter's cyclone rings arise from
discrete vortex self-organization (a cluster of storms).  Both mechanisms
deliver the system to $|\lam| \approx 1$, where the polygon lives---and
once there, eigenvalue branch disconnection and negative-temperature
statistics keep it.

The paper is organized as follows.  \Cref{sec:setup} introduces the
mathematical framework: the QGPV equation, M\"obius transformations, and
the loxodromic flow family.  \Cref{sec:theorems} proves the main results:
stationarity equivalence, branch disconnection, Thomson critical ratio,
amplitude formula, and energy--enstrophy sign structure.
\Cref{sec:verification} tests the framework against Saturn (Cassini) and
Jupiter (Juno) observations.  \Cref{sec:discussion} assembles the
five-step answer.

\section{Mathematical Setup}
\label{sec:setup}

\subsection{QGPV equation on the \texorpdfstring{$\beta$}{beta}-plane}

The barotropic quasi-geostrophic potential vorticity (QGPV) equation on the
$\beta$-plane reads
\begin{equation}\label{eq:qgpv}
  \frac{Dq}{Dt} = 0, \qquad q = \nabla^2\psi + \beta y,
\end{equation}
where $\psi$ is the stream function, $\beta = \dd f/\dd y$ is the
meridional gradient of the Coriolis parameter, and $D/Dt = \partial_t +
J(\psi, \cdot)$ is the material derivative.

\subsection{Log-polar coordinate transformation}

Under the conformal map $(x,y) \mapsto (\rho,\theta)$ defined by
$x + iy = e^{\rho + i\theta}$, the Laplacian transforms as
\begin{equation}\label{eq:laplacian-logpolar}
  \nabla^2 = e^{-2\rho}\!\left(
    \frac{\partial^2}{\partial\rho^2} +
    \frac{\partial^2}{\partial\theta^2}
  \right).
\end{equation}
Linearizing \eqref{eq:qgpv} about a zonal base state $U(\rho)$ and Fourier
expanding the perturbation $\psi' = \psihat_n(\rho)\,e^{in\theta}\,
e^{-i\omega t}$ yields the Rayleigh--Kuo eigenvalue ODE
\begin{equation}\label{eq:rayleigh-kuo}
  (U - c)\bigl(\psihat_n'' - n^2 \psihat_n\bigr)
  + (\beta - U'')\,\psihat_n = 0,
\end{equation}
where $c = \omega/n$ is the phase speed and primes denote $\dd/\dd\rho$.
For stationary waves ($c = 0$) this simplifies to
\begin{equation}\label{eq:stationary-ode}
  \psihat_n'' + \left[\frac{\beta - U''}{U} - n^2\right] \psihat_n = 0.
\end{equation}

\subsection{M\"obius transformations and the conjugacy theorem}
\label{sec:mobius-conjugacy}

A M\"obius (linear fractional) transformation
$f(z) = (az+b)/(cz+d)$ with $ad - bc \neq 0$ acts on the Riemann sphere
$\mathbb{C}_\infty = \mathbb{C} \cup \{\infty\}$.  Its classification is
determined by the \emph{multiplier} $\lam$, defined as follows.

\begin{theorem}[M\"obius conjugacy]\label{thm:conjugacy}
Let $f$ be a non-parabolic M\"obius transformation with distinct fixed
points $z_1, z_2 \in \mathbb{C}_\infty$.  Define the conjugating map
$g(z) = (z - z_1)/(z - z_2)$.  Then:
\begin{enumerate}
  \item $h = g \circ f \circ g^{-1}$ satisfies $h(w) = \lam w$, where
    $\lam = f'(z_1)$ is the derivative at the fixed point $z_1$.
  \item The multiplier satisfies $\lam = \mu_1/\mu_2$ where $\mu_1, \mu_2$
    are the eigenvalues of the associated matrix
    $\bigl[\begin{smallmatrix}a&b\\c&d\end{smallmatrix}\bigr]$.
  \item Writing $\lam = e^{\sigma + i\alpha}$ with
    $\sigma = \log|\lam|$, $\alpha = \arg(\lam)$:
    \begin{center}
    \begin{tabular}{lll}
    \toprule
    $\sigma \neq 0$, $\alpha \neq 0$ & loxodromic & logarithmic spiral orbits \\
    $\sigma = 0$, $\alpha \neq 0$    & elliptic   & circular orbits \\
    $\sigma \neq 0$, $\alpha = 0$    & hyperbolic & radial orbits \\
    \bottomrule
    \end{tabular}
    \end{center}
\end{enumerate}
\end{theorem}

\begin{proof}
For $c \neq 0$, the fixed points of $f$ satisfy $cz^2 + (d-a)z - b = 0$
with discriminant $\Delta = (a+d)^2 - 4(ad-bc) \neq 0$ (non-parabolic),
giving two distinct finite fixed points.  For $c = 0$, $f$ is affine with
one finite fixed point and $z_2 = \infty$; the conjugacy proceeds
identically with $g(z) = z - z_1$.
Since $g(z_1) = 0$ and $g(z_2) = \infty$, the map $h$ fixes both $0$ and
$\infty$.  The only M\"obius maps fixing both are $w \mapsto \lam w$.  By
the chain rule, $\lam = h'(0) = g'(z_1) \cdot f'(z_1) \cdot
(g^{-1})'(0) = f'(z_1)$.
\end{proof}

\begin{remark}[Fixed point character]\label{rem:fixed-points}
For loxodromic and hyperbolic transformations ($|\lam| \neq 1$), $z_1$
is repelling and $z_2$ is attracting (or vice versa, depending on the
convention $|\lam| \ge 1$).  For \emph{elliptic} transformations
($|\lam| = 1$), both fixed points are \emph{neutral}---centers of
rotation, neither attracting nor repelling.  The term ``repelling fixed
point'' used in the loxodromic context does not apply to the elliptic
case.
\end{remark}

\subsection{From circles to polygons: the \texorpdfstring{$\mathbb{Z}_N$}{Z\_N}
projection}
\label{sec:circle-to-polygon}

The conjugacy theorem establishes that elliptic ($\sigma = 0$) orbits are
\emph{circles}, not polygons.  The regular $N$-gon arises through a
separate geometric step.

\begin{proposition}[Polygonal meander from $\mathbb{Z}_N$ perturbation]
\label{prop:polygon}
Any axisymmetric jet perturbed by an $n = N$ Fourier mode produces a
jet boundary
\begin{equation}\label{eq:jet-boundary}
  \rho_{\mathrm{boundary}}(\theta) = \rhostar + \varepsilon\cos(N\theta + \phi)
\end{equation}
that is a meander with $N$-fold symmetry, approximating a regular $N$-gon
for small $\varepsilon$.
\end{proposition}

\begin{proof}
The total PV field with the $n = N$ Fourier perturbation is
$q(\rho, \theta) = q_0(\rho) + q_n(\rho)\cos(N\theta + \phi)$.
The jet boundary is the PV contour $q = q_B$ (constant).  Linearizing
around $\rho = \rhostar$ where $q_0(\rhostar) = q_B$:
\[
  q_0'(\rhostar)\,\delta\rho + q_n(\rhostar)\cos(N\theta + \phi) = 0,
\]
giving $\delta\rho(\theta) = -\frac{q_n(\rhostar)}{q_0'(\rhostar)}
\cos(N\theta + \phi)$, which is \eqref{eq:jet-boundary} with
$\varepsilon = |q_n(\rhostar)/q_0'(\rhostar)|$.

This meander has $N$ maxima at $\theta = 2\pi k/N$ ($k = 0, \ldots, N-1$),
all at the same radius $R(1 + \varepsilon)$, equally spaced in angle---the
vertices of a regular $N$-gon to leading order in $\varepsilon$.  (At
$O(\varepsilon^2)$, the chord lengths between adjacent vertices differ
because the meander curves between maxima; the polygon approximation is
valid for $\varepsilon \ll 1$.)
\end{proof}

\begin{remark}[Role of Proposition~\ref{prop:polygon}]
This result is a standard fact about Fourier perturbations and does not
require M\"obius theory.  Its role in the paper is to provide the
geometric bridge between the M\"obius conjugacy theorem (which gives
\emph{circles} at $\sigma = 0$) and the observed \emph{polygons}: the
polygon arises when the QGPV selects a specific azimuthal mode $n = N$
(\cref{sec:theorems}), which converts the circular base state into an
$N$-gon meander via this proposition.
\end{remark}

This distinction matters: the M\"obius orbit is a circle; the polygon
emerges from the \emph{interaction} between the circular orbit structure
and the discrete $\mathbb{Z}_N$ symmetry imposed by the wavenumber
selection (Rossby dispersion or vortex dynamics).  The two steps are:
\begin{enumerate}
  \item $\sigma = 0$ places the system on the circular-orbit family
    (M\"obius conjugacy theorem);
  \item $\alpha = 2\pi/N$ restricts this to the $\mathbb{Z}_N$-symmetric
    sector, producing the $N$-gon meander.
\end{enumerate}

\subsection{Loxodromic flow family and orbit structure}

The continuous one-parameter group generated by $h(w) = \lam w$ is
$w(t) = e^{(\sigma + i\alpha)t}\,w_0$.  In polar coordinates
$w = Re^{i\Theta}$:
\begin{equation}\label{eq:orbit-polar}
  R(t) = R_0\,e^{\sigma t}, \qquad \Theta(t) = \alpha t + \Theta_0.
\end{equation}
Eliminating $t$ (for $\alpha \neq 0$):
\begin{equation}\label{eq:spiral-equation}
  \ln R = \ln R_0 + \frac{\sigma}{\alpha}(\Theta - \Theta_0).
\end{equation}
This is a logarithmic spiral with constant winding angle
$\psi = \arctan(\alpha/\sigma)$.  At $\sigma = 0$: $R = R_0$ for all
$\Theta$ (circular orbit, $\psi = \pi/2$).  As $\sigma \to 0^+$: the
spiral opens continuously to a circle.  Both $R(\Theta; \sigma)$ and
$\psi(\sigma)$ are $C^\infty$ in $\sigma$ for fixed $\alpha \neq 0$.

\begin{remark}[Stream function and the periodic domain]
\label{rem:stream-function}
The full complex potential $w(z) = Az^s$ with $s = \sigma + i\alpha$ gives
\begin{equation}\label{eq:psi-full}
  \psi_{\mathrm{full}}(\rho,\theta)
  = |A|\,e^{\sigma\rho - \alpha\theta}
  \sin(\alpha\rho + \sigma\theta + \arg A).
\end{equation}
The factor $e^{-\alpha\theta}$ is not $2\pi$-periodic in $\theta$ for
$\alpha \neq 0$.  This function is properly defined on the
\emph{universal cover} of the punctured plane---the strip
$\{(\rho,\theta) : \rho \in \mathbb{R},\; \theta \in \mathbb{R}\}$
with the identification $z = e^{\rho + i\theta}$---rather than on
the periodic domain $\theta \in [0, 2\pi)$.

\emph{We do not use $\psi_{\mathrm{full}}$ as a periodic stream function.}
The QGPV analysis proceeds entirely through the Fourier-decomposed
eigenvalue equation \eqref{eq:rayleigh-kuo}, which governs the radial
eigenfunction $\psihat_n(\rho)$ for each azimuthal mode $n$ independently.
The full perturbation
\begin{equation}\label{eq:fourier-mode}
  \psi'(\rho,\theta,t) = \psihat_n(\rho)\,e^{in\theta}\,e^{-i\omega t}
\end{equation}
\emph{is} $2\pi$-periodic in $\theta$ and satisfies the linearized QGPV
by construction (it is the Fourier decomposition of the linearized
equation).

The complex potential $w(z) = Az^s$ enters only through the \emph{orbit
geometry} of the M\"obius transformation $f(z) = \lam z$
(\cref{sec:mobius-conjugacy})---specifically, the shape of the orbits
as a function of $\sigma$---and not as a stream function substituted
into the QGPV.  The bridge between the M\"obius geometry and the QGPV
is the stationarity theorem (Theorem~\ref{thm:stationarity}), which
operates on $\psihat_n(\rho)$, not on $\psi_{\mathrm{full}}$.
\end{remark}

\section{Main Results}
\label{sec:theorems}

Throughout this section, $\sigma$ denotes the M\"obius parameter
$\sigma = \log|\lam| = \mathrm{Re}(s)$ unless otherwise qualified.  The
jet half-width is denoted $\ell_J$ (not $\sigma_J$) to avoid confusion.

\subsection{Wavenumber selection (Saturn)}

\begin{proposition}[Stationary wavenumber]\label{prop:wavenumber}
For the barotropic Rossby dispersion relation $c = U - \beta/k^2$ with
$k = n/R_{\mathrm{hex}}$, the stationary condition $c = 0$ gives
\begin{equation}\label{eq:nstar}
  n^* = R_{\mathrm{hex}}\sqrt{\beta / U_{\max}}.
\end{equation}
For Saturn ($R_{\mathrm{hex}} = 5.1 \times 10^7$~m,
$\beta = 1.46 \times 10^{-12}$~m$^{-1}$\,s$^{-1}$,
$U_{\max} = 120$~m/s): $n^* = 5.62$, whose nearest integer is $n = 6$.
\end{proposition}

\begin{proof}
The Rossby dispersion relation for a barotropic wave with zonal
wavenumber $k$ and zero meridional wavenumber is $c = U - \beta/k^2$.
Setting $c = 0$ (stationary wave): $U = \beta/k^2$, giving
$k = \sqrt{\beta/U}$.  With $k = n/R_{\mathrm{hex}}$:
\[
  n = R_{\mathrm{hex}}\,k = R_{\mathrm{hex}}\sqrt{\beta/U}.
\]
Substituting Saturn's parameters:
\[
  n^* = 5.1 \times 10^7 \times \sqrt{\frac{1.46 \times 10^{-12}}{120}}
  = 5.1 \times 10^7 \times 1.103 \times 10^{-7} = 5.62.
\]
The nearest integer is $n = 6$.  The next candidates ($n = 5$: $U^* = 152$~m/s;
$n = 7$: $U^* = 77$~m/s) are substantially farther from the observed
$U_{\mathrm{obs}} = 120$~m/s.
\end{proof}

\subsection{Stationarity equivalence}

\begin{theorem}[Stationarity equivalence]\label{thm:stationarity}
Within the WKB approximation near the jet center, if the Rayleigh--Kuo
potential satisfies $V(\rhostar) < 0$ (a condition verified for Saturn's
jet parameters), then the stationary Rossby wave eigenfunction has
$\sigma = 0$, i.e., $|\lam| = 1$.
\end{theorem}

\begin{proof}
Near the jet center, the Rayleigh--Kuo potential
$V(\rho) = n^2 - (\beta - U'')/U$ is approximately constant over a
region of width $\sim \ell_J$.  Within this region, the WKB approximation
gives $\psihat_n(\rho) \approx C\,e^{s\rho}$ with $s^2 = V(\rhostar)$
and $s$ constant to leading order.

Since $V \in \mathbb{R}$, we have $s^2 \in \mathbb{R}$.  Writing
$s = \sigma + i\alpha_\rho$ (where $\alpha_\rho$ is the WKB radial
oscillation rate):
\[
  \mathrm{Im}(s^2) = 2\sigma\alpha_\rho = 0.
\]
If $V(\rhostar) < 0$, the solution is oscillatory: $s = \pm i\alpha_\rho$
with $\alpha_\rho = \sqrt{-V} > 0$.  Since $\alpha_\rho \neq 0$, we
conclude $\sigma = 0$.

\emph{Verification of $V < 0$:}  For Saturn's Gaussian jet at the jet
center ($\rho = \rhostar$), $(\beta - U''(0))/U(0) \approx
\ell_J^{-2} + \beta/U_{\max} \approx 400 + 10^{-16} \approx 400$
(dominated by the jet curvature), while $n^2 = 36$.  Therefore
$V(0) = 36 - 400 = -364 < 0$.  The condition holds by a wide margin.

With $\sigma = 0$, the real part gives the stationary jet speed:
$U^* = (\beta - U''(\rhostar)) / [(\alpha_\rho^2 + n^2)\,e^{-2\rhostar}]$.
For Saturn: $U^* \approx 105$~m/s, consistent with
$U_{\mathrm{obs}} = 120.1$~m/s \citep{Antunano2015}
(12\% discrepancy from the barotropic approximation).
\end{proof}

\begin{remark}[The identification $\alpha = 2\pi/n$]\label{rem:alpha}
The M\"obius azimuthal parameter $\alpha$ in $\lam = e^{\sigma + i\alpha}$
is \emph{identified} with $2\pi/n$ from the $\mathbb{Z}_N$ symmetry of
the polygon: an $N$-gon with vertices at $\theta_k = 2\pi k/N$ has a
discrete rotational symmetry of angle $2\pi/N$, which is by definition
$\arg(\lam)$ for the rotation $z \mapsto e^{i\,2\pi/N} z$.

This is not a derivation from the M\"obius framework---it is a restatement
of the definition of the multiplier for a linear rotation.  The M\"obius
framework contributes the \emph{continuous} family parameterized by
$\sigma$ (the spiral-to-circle transition); the \emph{discrete}
wavenumber $n$ comes from the QGPV (Proposition~\ref{prop:wavenumber}).

The WKB radial oscillation rate $\alpha_\rho$ satisfying
$\alpha_\rho^2 = |V(\rhostar)|$ is a separate quantity.  This paper
uses $\alpha$ for the M\"obius azimuthal parameter and $\alpha_\rho$ for
the WKB rate.
\end{remark}

\subsubsection{Branch disconnection within WKB}

\begin{proposition}[Branch disconnection]\label{prop:protection}
Within the WKB approximation, the condition $\sigma = 0$ is preserved
under perturbations of the jet speed $\delta U$.  The elliptic
($\sigma = 0$) and hyperbolic ($\alpha = 0$) branches of the eigenvalue
equation are disconnected.
\end{proposition}

\begin{proof}
The potential $V(\rhostar; \delta U)$ is real for all real $\delta U$.
The eigenvalue equation $s^2 = V$ requires $s^2 \in \mathbb{R}$.
Writing $s = i\alpha + \delta s$ with $\delta s = p + iq$:
\[
  \mathrm{Im}(s^2) = 2p(\alpha + q) = 0.
\]
This gives two branches:
\begin{enumerate}
  \item \textbf{Elliptic} ($p = 0$): $\sigma = 0$ for all $\delta U$.
  \item \textbf{Hyperbolic} ($q = -\alpha$): $\sigma = \sqrt{V + \alpha^2}$,
    purely radial flow.
\end{enumerate}
These branches meet only at $V = 0$ ($\delta U/U^* = \alpha^2/n^2
\approx 0.03$ for Saturn) and are not continuously connected for
$V \neq 0$.
\end{proof}

\begin{remark}[Scope of the protection]\label{rem:protection-scope}
Proposition~\ref{prop:protection} establishes branch disconnection
\emph{within the WKB approximation}, which is valid to leading order in
$\ell_J$.  Higher-order corrections to the eigenvalue equation could in
principle connect the branches; we have not excluded this.  The physical
persistence of Saturn's hexagon across four decades is consistent with
the branch disconnection but is not a consequence of
Proposition~\ref{prop:protection} alone.
\end{remark}

\subsection{The M\"obius parameter from the boundary value problem}

\begin{remark}[Multiple $\sigma$-like quantities]\label{rem:sigmas}
The M\"obius parameter $\sigma = \log|\lam|$ should be distinguished from:
$\sigma_{\mathrm{mode}} = \mathrm{Re}(s)$ in the QGPV eigenfunction
(which is zero within WKB); $\ell_J$, the jet half-width in log-polar
units; and $\sigma_{\mathrm{geom}}$, the observational polygon distortion
measure defined in \cref{sec:sigma-geom}.  Only $\sigma$ (the M\"obius
parameter) and $\sigma_{\mathrm{geom}}$ (the observable) appear in the
cross-planetary comparison.  They measure different things:
$\sigma$ characterizes the mean flow geometry (spiral winding);
$\sigma_{\mathrm{geom}}$ measures deviation from a perfect $N$-gon.
\end{remark}

The M\"obius parameter $\sigma$ is determined from the three-region
boundary value problem through a two-step chain:
\begin{enumerate}
  \item The BVP matching (below, Theorem~\ref{thm:matching}) determines
    the hexagon amplitude $\varepsilon$.
  \item The hexagonal boundary geometry gives the velocity ratio
    $|u_\rho/u_\theta|_{\mathrm{rms}} = n\varepsilon/\sqrt{2}$, and
    $\sigma = \alpha \cdot |u_\rho/u_\theta|$.
\end{enumerate}

\subsection{Thomson vortex ring stability}

\begin{proposition}[Critical central vortex strength]\label{prop:critical-ratio}
For the $N = 6$ Thomson ring on the unit circle with central vortex
$\kappa_0$, the characteristic polynomial of the $12 \times 12$
Hamiltonian stability matrix (constructed symbolically in Mathematica) is
\begin{equation}\label{eq:char-poly}
  p(\lam) \propto \lam^2 \cdot
  \underbrace{\bigl(9 + 36\kappa_0/\kappa + 16\pi^2\lam^2\bigr)}_{\text{Factor A}} \cdot
  \underbrace{Q(\lam^2, \kappa_0/\kappa)^2}_{\text{Factor B (double)}},
\end{equation}
where $Q$ is a quartic in $\lam^2$ with perfect-square discriminant
$(9 - 12\kappa_0/\kappa)^2$.  The first instability arises from Factor~A:
\begin{equation}\label{eq:critical-ratio}
  \boxed{\left(\frac{\kappa_0}{\kappa}\right)_{\!\mathrm{crit}} = -\frac{1}{4}.}
\end{equation}
\end{proposition}

\begin{proof}
The Jacobian of the Hamiltonian point vortex equations in the co-rotating
frame, including ring--ring and ring--center interactions, is constructed
symbolically.  Its characteristic polynomial, computed by Mathematica's
\texttt{Det} function and verified numerically at multiple $\kappa_0$
values, factors as \eqref{eq:char-poly}.  The stability threshold from
Factor~A is $9 + 36\kappa_0/\kappa = 0$, giving
$\kappa_0/\kappa = -1/4$.
\end{proof}

Remarkably, $\kappa_{\mathrm{crit}} = -1/4$ is specific to $N = 6$: for
$N = 4$ the critical ratio is $-1/2$ (proven algebraically via
characteristic polynomial), and numerical computation gives
$\kappa_{\mathrm{crit}} = -1/2$ for $N = 3$ and $N = 5$ as well
(to within $10^{-3}$; algebraic proof for these cases is obstructed by
irrational vertex coordinates).  The hexagonal ring thus has enhanced
sensitivity to anticyclonic central perturbations compared to smaller
rings.

\paragraph{Spectral non-equivalence.}
The Thomson stability matrix and the Rayleigh--Kuo operator share
$\mathbb{Z}_6$ eigenvectors (discrete Fourier modes) but not eigenvalues.
Galerkin projection gives eigenvalue ratios varying from $2.74$ (mode~1)
to $0.14$ (mode~3).  The operators are \emph{not} spectrally equivalent;
they converge at $|\lam| = 1$ independently.

\subsection{Matched asymptotic amplitude (Saturn)}

\begin{theorem}[Hexagon amplitude]\label{thm:matching}
The matched asymptotic expansion connecting the inner loxodromic flow
to the outer Rossby wave through the PV step at $\rhostar$ predicts the
hexagon amplitude
\begin{equation}\label{eq:amplitude}
  \varepsilon = C \cdot \frac{\delta U}{U^*} \cdot
  \exp\!\bigl(-n\,|\rhostar|\bigr),
\end{equation}
where $\delta U = U_{\mathrm{obs}} - U^*$ and $C$ is determined by the
Fourier projection of the PV gradient onto the eigenfunction envelope:
\begin{equation}\label{eq:overlap-integral}
  C = \frac{1}{U^* n}\int_0^\infty
  \bigl(\beta - U''(\rho)\bigr)\,e^{-n\rho}\,\dd\rho.
\end{equation}
For the Gaussian jet $U(\rho) = U_{\max}\exp(-\rho^2/2\ell_J^2)$:
\begin{equation}\label{eq:matching-constant}
  C = \frac{U_{\max}}{U^*} - \frac{n\,\ell_J}{U^*}
      \sqrt{\frac{\pi}{2}}\,e^{n^2\ell_J^2/2}\,
      \mathrm{Erfc}\!\left(\frac{n\,\ell_J}{\sqrt{2}}\right)
      + \frac{\beta}{n\,U^*}.
\end{equation}
For Saturn ($n = 6$, $\ell_J = 0.05$, $U_{\max} = 120$~m/s,
$U^* = 105$~m/s): $C = 0.797$, giving $\varepsilon = 0.112$.
\end{theorem}

This matches the Cassini-observed meander amplitude
$\varepsilon_{\mathrm{obs}} = 0.112$, under a Gaussian jet approximation
(9.8\% RMS fit to the actual wind profile; see \cref{sec:discussion} for
the propagated uncertainty).

\subsection{Energy--enstrophy sign structure}
\label{sec:enstrophy}

\begin{proposition}[Energy maximum at the polygon]\label{prop:energy-max}
For $N \ge 3$ point vortices on a ring, consider the spiral deformation
\begin{equation}\label{eq:spiral-deformation}
  z_k(\sigma) = R\,e^{\sigma k/N}\,e^{2\pi ik/N},
  \qquad k = 0, \ldots, N-1,
\end{equation}
which at $\sigma = 0$ gives the regular $N$-gon and for $\sigma > 0$
stretches vertex $k$ outward by the factor $e^{\sigma k/N}$, spiraling
the configuration.  The Thomson interaction energy
\begin{equation}\label{eq:H-sigma}
  H(\sigma) = -\sum_{j<k} \ln\bigl|z_j(\sigma) - z_k(\sigma)\bigr|
\end{equation}
satisfies $H''(0) < 0$.
\end{proposition}

\begin{remark}[Nature of the deformation]
The deformation \eqref{eq:spiral-deformation} is \emph{not} the uniform
M\"obius action $z_k \mapsto \lam z_k$ (which scales and rotates the
entire polygon rigidly, leaving $H$ unchanged).  It is a
\emph{symmetry-breaking} deformation that assigns each vertex a different
radial stretch, converting the $N$-gon into a logarithmic spiral.  This
direction in configuration space is physically relevant: it is the
deformation produced by a loxodromic flow $w(z) = Az^s$ acting on fluid
parcels at different azimuthal positions over time.  However, we do not
claim this is the \emph{unique} relevant deformation direction; $H''(0)$
could have different signs along other directions in the $N$-dimensional
configuration space.  The result is that the polygon is an energy maximum
\emph{along the spiral deformation direction}, not necessarily globally.
\end{remark}

\begin{proof}
At $\sigma = 0$, the pairwise squared distance is
$|z_j - z_k|^2 = 4R^2\sin^2(\pi m/N)$ where $m = k - j$.
Define $F(u, m) = \ln[1 - 2\cos(2\pi m/N)\,e^{um} + e^{2um}]$ where
$u = \sigma/N$, so that $\ln|z_j - z_k|^2 = 2uj + \ln R^2 + F(u,m)$.
Differentiating $F$ twice:
\[
  F''(0, m) = \frac{m^2}{1 - \cos(2\pi m/N)}
  = \frac{m^2}{2\sin^2(\pi m/N)}.
\]
Since $\dd^2(2uj)/\dd u^2 = 0$, the only contribution to $H''$ is from $F$.
Each pair $(j,k)$ with separation $m$ contributes $F''(0,m)$, and there
are $(N - m)$ such pairs for $m = 1, \ldots, N-1$.  Converting
$\dd/\dd\sigma = (1/N)\,\dd/\dd u$:
\begin{equation}\label{eq:Hpp-formula}
  H''(0) = -\frac{1}{4N^2}\sum_{m=1}^{N-1}
  \frac{(N-m)\,m^2}{\sin^2(\pi m/N)}.
\end{equation}
Every factor in the sum is strictly positive: $(N-m) > 0$, $m^2 > 0$,
$\sin^2(\pi m/N) > 0$ for $m = 1, \ldots, N-1$.  Therefore the sum is
positive and $H''(0) < 0$.
\end{proof}

The regularized enstrophy $Z(\sigma) = \sum_{j<k}|z_j - z_k|^{-2}$
satisfies $Z''(0) > 0$ (verified computationally for $N = 3, \ldots, 9$).
The ratio $\mu = -Z''/H'' > 0$ for all $N$, corresponding to negative
temperature in the Onsager sense.

\begin{remark}[Scope of the energy maximum]
The curvatures $H''(0)$ and $Z''(0)$ are computed along the specific
one-parameter loxodromic deformation defined above.  In the full
$N$-dimensional configuration space, perturbations exist that decrease
$Z$ (random tests show 62--78\% of perturbations increase $Z$, not
100\%).  The energy maximum is \emph{local} and
\emph{direction-dependent}, not global.

While $H''(0) < 0$ is proven for all $N \ge 3$
(Proposition~\ref{prop:energy-max}), the enstrophy result $Z''(0) > 0$
remains computational.
\end{remark}

\section{Observational Verification}
\label{sec:verification}

The framework makes quantitative predictions for both Saturn and Jupiter.
We test them against Cassini and Juno observations, treating the two
planets as independent case studies of the same geometric principle.

\subsection{A generalized observable: \texorpdfstring{$\sigma_{\mathrm{geom}}$}{sigma-geom}}
\label{sec:sigma-geom}

To compare polygonal patterns across planets and dynamical mechanisms,
we define a universal observable.  Given $N$ features (cyclone centers,
jet meander vertices) at positions $z_k$ in the complex plane (via
stereographic projection from the pole), the geometric M\"obius parameter
is
\begin{equation}\label{eq:sigma-geom}
  \sigma_{\mathrm{geom}} = \frac{1}{N}\sum_{k=1}^{N}
  \bigl|\!\ln|z_k / z_k^{\mathrm{ideal}}|\bigr|,
\end{equation}
where $z_k^{\mathrm{ideal}}$ are the vertices of the best-fit regular
$N$-gon (optimized over rotation angle).  By construction,
$\sigma_{\mathrm{geom}} = 0$ for a perfect polygon ($|\lam| = 1$) and
$\sigma_{\mathrm{geom}} > 0$ for any distortion.  The observable is
rotation-invariant, scale-invariant, and independent of $N$.

\subsection{Saturn: the hexagonal jet meander}

\subsubsection{Three-epoch data}

Saturn's hexagon has been observed across three epochs spanning
three decades (\cref{tab:epochs}).

\begin{table}[htbp]
\centering
\caption{Saturn hexagon drift rates across three observational epochs.
All angular velocities relative to System~III.}
\label{tab:epochs}
\begin{tabular}{lcccl}
\toprule
Epoch & $\omega_{\mathrm{hex}}$ ($^\circ$/day) & $\omega_{\mathrm{NPS}}$ ($^\circ$/day) & NPS & Source \\
\midrule
1980--81  & $-0.060 \pm 0.014$ & $-0.044 \pm 0.010$ & Yes   & \citet{Godfrey1988} \\
1990--91  & $-0.001 \pm 0.007$ & ---                  & Weak  & \citet{SanchezLavega1993} \\
2008--14  & $-0.036 \pm 0.010$ & ---                  & No    & \citet{SanchezLavega2014} \\
\bottomrule
\end{tabular}
\end{table}

\subsubsection{Wavenumber and amplitude}

The Cassini/ISS zonal wind profile \citep{Antunano2015}, digitized near
the hexagonal jet ($70$--$82^\circ$\,N), gives a peak speed
$U_{\mathrm{obs}} = 120.1$~m/s at $76.6^\circ$\,N with half-width
$2.86^\circ$.  The Rossby stationarity condition $U^* = \beta/k^2$ with
$k = 6/R_{\mathrm{hex}}$ predicts $U^* = 105$~m/s, selecting $n^* = 5.62
\to 6$ as the stationary wavenumber (12\% discrepancy attributable to the
barotropic approximation).

The hexagon meander amplitude, measured from the vertex latitude span
($75^\circ$--$78^\circ$\,N), is $\varepsilon_{\mathrm{obs}} = 0.112$.
The Theorem~5 prediction with the derived matching constant $C = 0.797$
gives $\varepsilon_{\mathrm{pred}} = 0.112$.  The numerical coincidence
with $\varepsilon_{\mathrm{obs}} = 0.112$ should not be overstated:
propagating the 9.8\% Gaussian profile approximation error gives
$\varepsilon_{\mathrm{pred}} = 0.11 \pm 0.11$, so the match is within
one standard deviation of the model uncertainty, not sub-percent.

\subsubsection{NPS coupling and near-resonance}

The North Polar Spot (NPS), a compact anticyclonic vortex near the
hexagon, was prominent during 1980--81 but absent by 2008--14.
The linear coupling model $\omega_{\mathrm{hex}} = \omega_0 +
\gamma\,\omega_{\mathrm{NPS}}$ gives $\gamma = 0.54 \pm 0.41$
(two parameters fitted to two data points with no residual degrees of
freedom; the 1990--91 epoch is not used because $\omega_{\mathrm{NPS}}$
was not measured).  The implied near-resonance period is
\[
  T_{\mathrm{rel}} = \frac{360^\circ}{|\omega_{\mathrm{hex}} -
  \omega_{\mathrm{NPS}}|} = 62 \pm 68\;\text{years},
\]
where the large uncertainty reflects the propagated error bars on both
drift rates.  The multi-decadal timescale is \emph{suggestive} of an
NPS--hexagon coupling but cannot be considered well-constrained by the
available data.

\subsection{Jupiter: the cyclone rings}

Jupiter provides the critical universality test.  Its polar cyclone
rings arise through vortex self-organization, not Rossby wave stationarity
---the Rossby prediction gives $n^* \approx 0.3$ for both poles,
completely failing to match the observed $N = 8$ (north) or $N = 5$
(south).  If the M\"obius framework applies only to Rossby-mediated
polygons, it should fail for Jupiter.  It does not.

\subsubsection{Juno cyclone positions}

Cyclone positions digitized from \citet{Adriani2018} are projected
stereographically from the pole and measured via
$\sigma_{\mathrm{geom}}$ (\cref{eq:sigma-geom}).

\subsubsection{Thomson stability}

The north polar ring ($N = 8$) is unstable without a central vortex
(Havelock: $N \ge 8$ unstable), but Jupiter has a strong central cyclone
at each pole that provides stabilization.  The south polar ring ($N = 5$)
is stable without a center (Havelock: $N \le 7$), with Thomson critical
ratio $\kappa_{\mathrm{crit}} = -1/2$.

\subsection{Cross-planetary comparison}

\paragraph{Quantitative criterion.}
We say a configuration is ``near $|\lam| = 1$'' if its M\"obius parameter
satisfies $\sigma < 1$ (equivalently $r = e^\sigma < e \approx 2.72$).
This threshold is motivated physically: $\sigma < 1$ means the spiral
pitch angle $\psi = \arctan(\alpha/\sigma)$ exceeds $45^\circ$---the
orbit is closer to a circle than to a radial ray.  We further distinguish
\emph{close} ($\sigma < 0.1$, $r < 1.11$) from \emph{moderate}
($0.1 < \sigma < 1$).

For Saturn, we note that two different observables give different values:
\begin{itemize}
  \item $\sigma_{\mathrm{flow}} \approx 0.10$ (from the dynamical BVP
    chain, \cref{sec:theorems})---\emph{close} to $|\lam| = 1$.
  \item $\sigma_{\mathrm{geom}} = 0.50$ (from the vertex distortion
    measure, \cref{sec:sigma-geom})---\emph{moderate}.
\end{itemize}
The difference arises because $\sigma_{\mathrm{geom}}$ measures the
full geometric distortion of the hexagonal meander (amplitude
$\varepsilon = 0.112$), while $\sigma_{\mathrm{flow}}$ measures the
Rossby wave's departure from exact stationarity ($\delta U/U^* = 0.14$).
Both place Saturn within the $\sigma < 1$ criterion; Jupiter sits in
the ``close'' regime by either measure.

\begin{table}[htbp]
\centering
\caption{All three planetary configurations tested against the framework.}
\label{tab:cross-planet}
\begin{tabular}{lccccl}
\toprule
System & $N$ & $\sigma_{\mathrm{geom}}$ & $\sigma_{\mathrm{flow}}$ & Regime & Entry mechanism \\
\midrule
Saturn hexagon    & 6 & 0.50 & 0.10 & Close (flow) / Moderate (geom) & Rossby \\
Jupiter north     & 8 & 0.06 & --- & Close & Thomson + center \\
Jupiter south     & 5 & 0.05 & --- & Close & Thomson \\
\bottomrule
\end{tabular}
\end{table}

Jupiter's cyclone rings ($\sigma_{\mathrm{geom}} < 0.07$) sit close to
$|\lam| = 1$.  Saturn's hexagon ($\sigma_{\mathrm{geom}} = 0.50$,
$r = 1.65$) departs substantially further.  This difference is physically
expected: $\sigma_{\mathrm{geom}}$ measures the \emph{geometric}
distortion of the observed feature from a perfect $N$-gon, and Saturn's
hexagon is a jet meander with amplitude $\varepsilon = 0.112$---a
large-scale deformation of a circular jet that necessarily distorts the
polygon.  Jupiter's cyclone rings, by contrast, are discrete point-like
vortices whose positions can closely approximate a regular polygon even
when the surrounding flow is not circular.  The two cases probe different
regimes of the same framework: Saturn tests the Rossby-mediated
\emph{wave} regime ($\sigma_{\mathrm{flow}} \approx 0.10$ from the
dynamical chain, closer to $|\lam| = 1$ than $\sigma_{\mathrm{geom}}$
suggests); Jupiter tests the vortex-cluster regime where
$\sigma_{\mathrm{geom}}$ directly measures the ring geometry.

A quantitative criterion for ``near $|\lam| = 1$'' depends on the
observable used: for the dynamical M\"obius parameter
$\sigma_{\mathrm{flow}}$, Saturn sits at $\sim 0.10$ (close); for the
geometric observable $\sigma_{\mathrm{geom}}$, Saturn sits at $0.50$
(moderate).  We report both and do not claim tighter agreement than the
data support.

The Rossby mechanism (Theorem~3) fails entirely for Jupiter.  The Thomson
mechanism provides the alternative entry point.  Different dynamics,
same geometric destination.

\begin{table}[htbp]
\centering
\caption{Summary of verified predictions.}
\label{tab:verification}
\begin{tabular}{clccl}
\toprule
\# & Prediction & Value & Observed & Status \\
\midrule
1 & Saturn wavenumber $n^*$ & $5.62$ & $6$ & Correct integer \\
2 & Saturn amplitude $\varepsilon$ & $0.11 \pm 0.11$ & $0.112$ (Cassini) & Consistent \\
3 & Stationary jet speed $U^*$ & $105$~m/s & $120$~m/s & $12\%$ (barotropic) \\
4 & Thomson $\kappa_{\mathrm{crit}}$ ($N=6$) & $-1/4$ & --- & Exact algebraic \\
5 & Branch disconnection & $\sigma = 0$ (within WKB) & 40-year persistence & Proven (WKB) \\
6 & Jupiter $\sigma_{\mathrm{geom}}$ & $< 0.07$ & Near-polygonal & Confirmed \\
\bottomrule
\end{tabular}
\end{table}

\section{Discussion}
\label{sec:discussion}

\subsection{Assembling the answer}

We can now answer the question posed in the title.

Rotating fluids make polygons because the regular $N$-gon occupies a
distinguished position in the M\"obius parameter space $\Cstar$: it sits
at $|\lam| = 1$, where spiral orbits become closed polygonal orbits
(Step~1, the conjugacy theorem).  This position is an energy maximum
among ring configurations of $N$ point vortices (Step~3, $H''(0) < 0$
for all $N \ge 3$).  In two-dimensional turbulence at negative
temperature---the regime of large-scale vortex condensation
\citep{Onsager1949}---systems are driven toward energy maxima (Step~4).
The regular polygon is therefore the most probable macrostate.
Atmospheric turbulence provides the thermostat that enables this
statistical relaxation (Step~5).  Once the system arrives at
$|\lam| \approx 1$, eigenvalue branch disconnection (within the WKB
approximation) prevents smooth perturbations of the jet speed from
moving the system to the hyperbolic branch (Step~2).

\emph{That} a polygon forms is universal---a consequence of geometry
(the existence of the $|\lam| = 1$ locus), statistics (negative
temperature selects energy maxima), and eigenvalue structure (the polygon
branch is disconnected from the radial branch within WKB).

\emph{Which} polygon forms is planet-specific.  Saturn's Rossby
dispersion relation gives $n^* = 5.62 \to 6$; Jupiter's vortex dynamics
selects $N = 5$ at the south pole and $N = 8$ at the north.  Different
mechanisms arrive at $|\lam| \approx 1$ through different doors---Rossby
stationarity for Saturn, Thomson stability for Jupiter---but they arrive
at the same geometric destination.

\subsection{What this framework does and does not explain}

\paragraph{Explained.}
Why polygons persist (branch disconnection within WKB).  Why the polygon is
statistically favored (Onsager negative temperature).  Why Saturn's
hexagon has six sides (Rossby wavenumber selection).  Why the amplitude
is $\varepsilon = 0.112$ (matched asymptotic expansion with derived $C$).
Why the hexagon--NPS interaction has a 62-year cycle (near-resonance).
Why Jupiter's cyclone rings are nearly perfect polygons
($\sigma_{\mathrm{geom}} < 0.07$).

\paragraph{Not explained.}
Why Jupiter has $N = 8$ at one pole and $N = 5$ at the other (this
requires detailed vortex dynamics modeling beyond the scope of this
framework).  Why the atmospheric thermostat operates at negative
temperature (this is observed but not derived from first principles).
Why the enstrophy is minimized along the loxodromic direction
($Z''(0) > 0$ is verified computationally but not proven analytically).

\subsection{The enstrophy--energy sign structure}
\label{sec:enstrophy-discussion}

The energy--enstrophy structure at $\sigma = 0$ deserves explicit
discussion because it connects the geometric framework to Onsager's
statistical mechanics.

Along the loxodromic deformation direction in $\Cstar$, the regular
$N$-gon has
\begin{equation}\label{eq:sign-structure}
  H''(0) < 0 \quad (\text{energy maximum}), \qquad
  Z''(0) > 0 \quad (\text{enstrophy minimum}),
\end{equation}
where $H''(0) < 0$ is proven for all $N \ge 3$
(Proposition~\ref{prop:energy-max}) and $Z''(0) > 0$ is verified
computationally for $N = 3, \ldots, 9$.  The ratio $\mu = -Z''/H'' > 0$
for all $N$, placing the polygon at the \emph{negative-temperature}
state: maximum energy, minimum enstrophy.

The enstrophy minimum is \emph{local}, restricted to the one-parameter
loxodromic deformation family.  In the full $N$-dimensional
configuration space, perturbations exist that decrease $Z$ (random tests
show 62--78\% of perturbations increase $Z$, not 100\%).  The polygon is
not the global enstrophy minimizer, but it occupies the distinguished
position along the physically relevant deformation direction.

The physical picture: once a dynamical mechanism delivers the system to
$\sigma \approx 0$, the energy maximum provides no spontaneous collapse
toward a dispersed state (that would require decreasing $H$, violating
energy conservation in the large-scale dynamics).  The enstrophy minimum
provides no force away from the polygon along the loxodromic direction.
Branch disconnection (within WKB) provides no perturbative exit from
$\sigma = 0$ at leading order.  The polygon is held in place by the
convergence of conservation laws, variational structure, and eigenvalue
branch structure.

\subsection{Gaps in the five-step argument}

Three logical gaps connect the five steps and should be stated explicitly.

\paragraph{The QGPV--vortex bridge.}
Steps 1--2 operate on the QGPV (a PDE for a continuous fluid); Steps 3--5
operate on the Thomson point vortex system (discrete vortices with
logarithmic interaction).  The universality claim requires that these two
descriptions apply to the same physical system.  For Jupiter's cyclone
rings (discrete storms), the point vortex model applies directly.  For
Saturn's hexagonal jet meander (a continuous flow), the identification
of the jet with a vortex ring is a modeling assumption, not a derivation.
The bridge between the two descriptions---when and why a continuous jet
can be represented by $N$ discrete vortices---is part of the broader
vortex-dynamics literature \citep{Aref2003} but is not established within
this paper.

\paragraph{Onsager applicability.}
Step 4 invokes Onsager's (1949) negative-temperature condensation for
point vortices.  Applying this to planetary atmospheres requires three
assumptions not proven here: (a) the large-scale flow is well-described
by a point vortex system, (b) the system is in statistical equilibrium
at negative temperature, and (c) the energy maximum of the point vortex
ring coincides with the polygonal state of the continuous flow.  These
are physically reasonable for Jupiter's cyclone rings (discrete vortices
in approximate statistical equilibrium) but less directly applicable to
Saturn's jet meander.

\paragraph{Amplitude uncertainty.}
The $\varepsilon_{\mathrm{pred}} = 0.112$ matching the Cassini
observation to $0.4\%$ overstates the precision.  The matching constant
$C = 0.797$ is computed from a Gaussian jet approximation with 9.8\%
RMS fit error.  Propagating this profile uncertainty through the
formula gives $\varepsilon_{\mathrm{pred}} = 0.11 \pm 0.11$---the
true uncertainty is of order 100\%, not $0.4\%$.  The match is within
one standard deviation of the Gaussian approximation, which is
consistent but not sharp.  Furthermore, $U^*$ appears in both the
numerator ($\delta U = U_{\mathrm{obs}} - U^*$) and denominator of
$\varepsilon = C \cdot \delta U / U^*$, making the prediction
sensitive to the \emph{internal consistency} of the Gaussian
approximation rather than to the actual wind profile.

\subsection{Limitations}

\begin{itemize}
  \item \textbf{Barotropic approximation.}  The 12\% discrepancy between
    $U^*$ and $U_{\mathrm{obs}}$ for Saturn reflects the neglect of
    vertical structure.

  \item \textbf{Digitized data.}  Cassini wind profiles and Juno cyclone
    positions are digitized from published figures.  Quantitative matches
    should be verified against PDS archival data.

  \item \textbf{Limited sample.}  Three planetary configurations from two
    planets.  Uranus and Neptune observations would strengthen or falsify
    the universality claim.

  \item \textbf{Statistical, not deterministic.}  The drive toward the
    polygon is statistical (most probable macrostate), not deterministic
    (Hamiltonian trajectory).  The atmospheric thermostat is observed but
    not derived from first principles.

  \item \textbf{Local enstrophy minimum.}  The sign structure
    \eqref{eq:sign-structure} holds along the loxodromic direction, not
    globally in configuration space.
\end{itemize}

\subsection{Future directions}

\begin{enumerate}
  \item Verify all predictions against PDS (Cassini) and JIRAM (Juno)
    archival data.
  \item Compute the Onsager partition function on $\Cstar$ to confirm
    that the negative-temperature equilibrium peaks at $\sigma = 0$.
  \item Extend to baroclinic QGPV to resolve the 12\% $U^*$ discrepancy.
  \item Apply the framework to Uranus and Neptune when high-resolution
    polar data become available.
  \item Derive the $N$-selection mechanism for Jupiter: does the central
    vortex strength $\kappa_0/\kappa$ determine $N$ through the Thomson
    critical ratio hierarchy?
  \item Test on rotating-tank experiments \citep{BarbosaAguiar2010}, where
    the thermostat (viscous dissipation) is known exactly.
\end{enumerate}


\bibliographystyle{plainnat}
\bibliography{refs}
\end{document}
